\chapter{Galois Cohomology and Field Arithmetic}
\label{ch:galois-cohomology}

Throughout, \(\operatorname{Gal}(L/k)\) for a possibly infinite Galois
extension has the Krull topology. Fix a separable closure \(k^{\mathrm s}\)
and write \(G_k=\operatorname{Gal}(k^{\mathrm s}/k)\). Superscript
\({}^{+}\) denotes an additive group. The coefficient groups and matrix
groups attached to a field in this chapter have the discrete topology.
We use the continuous inhomogeneous complex of Result GC01 in
\emph{Topics in Group Theory} \cite[PDF p.~24]{mg-groups}. In degree one,
with multiplicative notation, \(c(gh)=c(g)\,g(c(h))\); for abelian
coefficients the coboundary attached to \(b\) is \(g\mapsto b^{-1}g(b)\).
For a possibly nonabelian coefficient group the equivalence is
\[
 c'(g)=b^{-1}c(g)\,g(b).
\]
The resulting \(H^1\) is a pointed set in this case, with the identity
cocycle as base point.

\section{Finite layers and Hilbert's Theorem 90}

\begin{proposition}[Finite Galois layers and fixed coefficients]
\label{res:GH01}
Let \(L/k\) be Galois and put \(G=\operatorname{Gal}(L/k)\).
For each finite Galois intermediate field \(E\) put
\(U_E=\operatorname{Gal}(L/E)\) and
\(G_E=\operatorname{Gal}(E/k)\). For a discrete continuous abelian
\(G\)-module \(M\) there are identifications
\[
 H^q_{\mathrm{cts}}(G,M)
   \ \cong\ \varinjlim_E H^q(G_E,M^{U_E})\qquad(q\geq 0).
\]
They use inflation along restriction of field automorphisms and inclusion
of fixed coefficients. For a discrete group \(B\) with continuous
\(G\)-action, the analogous identification holds for the pointed sets
\(H^1_{\mathrm{cts}}(G,B)\).
\end{proposition}
\begin{proof}
Results F04, T02, T06 and C02 of \emph{Infinite Galois Theory}
\cite[PDF pp.~5, 9--10, 12]{mg-aft} give the directed finite layers,
the open normal basis \(U_E\), the quotient \(G/U_E=G_E\), and
\(L^{U_E}=E\). The abelian assertion is GC02 of \emph{Topics in Group
Theory} \cite[PDF p.~24]{mg-groups} applied to this cofinal basis.

For a nonabelian cocycle, compactness gives a finite image, and
Result PF03(4) of \emph{Topics in Group Theory}
\cite[PDF p.~16]{mg-groups} factors its underlying continuous map
through a finite quotient. Enlarge the layer \(E\) so that \(U_E\)
fixes all image elements, by intersecting their open stabilizers
with the first quotient kernel and choosing a basis subgroup.
Then the cocycle identity descends by surjectivity of restriction.
If two cocycles are related by some \(b\in B\), enlarge the common
layer once more so that \(U_E\) fixes \(b\). Their equivalence now
holds at that stage. These two finite-layer assertions give
surjectivity and injectivity of the pointed-set comparison.
\end{proof}

\begin{lemma}[A trace-one element]
\label{res:GH02}
If \(E/k\) is finite Galois, there exists \(t\in E\) with
\(\sum_{\sigma\in\operatorname{Gal}(E/k)}\sigma(t)=1\).
\end{lemma}
\begin{proof}
The independence lemma for distinct field embeddings in D07,
Lemma 16.6 \cite[PDF p.~125]{mg-d07}, makes the sum of the
automorphisms in the finite Galois group a nonzero function.
It is \(k\)-linear and its values are fixed by the Galois group, hence
are in \(k\) by the finite fixed-field theorem (Result F01 of
\emph{Infinite Galois Theory} \cite[PDF p.~4]{mg-aft}).
Divide an element of nonzero trace by its trace.
\end{proof}

\begin{theorem}[Matrix Hilbert 90]
\label{res:GH03}
For \(n\geq 1\) and a Galois extension \(L/k\), with entrywise action,
\[
 H^1_{\mathrm{cts}}\bigl(\operatorname{Gal}(L/k),
                         \operatorname{GL}_n(L)\bigr)=\{*\}.
\]
In particular \(H^1_{\mathrm{cts}}(\operatorname{Gal}(L/k),L^\times)=0\).
\end{theorem}
\begin{proof}
For a finite extension \(E/k\) and a cocycle with values in
\(\operatorname{GL}_n(E)\), Lemma~\ref{res:PV18} gives an invertible
matrix \(b\) with \(A_\sigma=b^{-1}\sigma(b)\) in any characteristic.
The stated equivalence applied to the identity cocycle shows
that this cocycle represents the base point.
In the infinite case \(\operatorname{GL}_n(L)^{U_E}=
\operatorname{GL}_n(E)\): the fixed entries lie in \(E\), and a
determinant nonzero in \(L\) is nonzero in \(E\).
Proposition~\ref{res:GH01} now reduces the assertion to the finite case.
\end{proof}

\begin{proposition}[Evaluation at a generator in degree one]
\label{res:GH04}
For \(C=\langle s\rangle\) of finite order \(d\) acting on an abelian
group \(A\), evaluation at \(s\) gives an isomorphism
\[
 H^1(C,A)\ \cong\ \ker(N:A\to A)/(s-1)A,\qquad
 N=1+s+\cdots+s^{d-1}.
\]
\end{proposition}
\begin{proof}
If \(d=1\), both sides are zero. Otherwise a cocycle with value \(a\)
at \(s\) has value \(\sum_{j=0}^{i-1}s^j a\) at \(s^i\) for
\(0\leq i<d\). At \(i=d\) the cocycle identity gives \(Na=0\);
conversely this condition makes the sum rule consistent under
reduction modulo \(d\) and implies the cocycle identity. A
coboundary has value \((s-1)b\) at \(s\), so quotienting gives
the stated isomorphism.
\end{proof}

Thus the older multiplicative and additive forms of cyclic
Hilbert 90 are statements that these particular quotients vanish.
D07 proves the norm-one and trace-zero formulas in Results
\texttt{K07.20.hilbert90\_multiplicative} and
\texttt{K07.20.hilbert90\_additive} \cite[PDF pp.~171--172]{mg-d07}.

\section{Kummer theory}

\begin{proposition}[Exact coefficient sequences for continuous cochains]
\label{res:GH05}
An exact sequence \(0\to A\to B\to C\to0\) of discrete abelian
continuous \(G\)-modules, for \(G\) profinite, gives the long exact
sequence in continuous cohomology. The connecting arrow sends
\([c]\) to \([d\widetilde c]\), for a continuous cochain lift
\(\widetilde c\) to \(B\).
\end{proposition}
\begin{proof}
The induced cochain map to \(C\) is surjective: a continuous map from
the compact space \(G^n\) to a discrete group takes finitely many
values, and the choice of one lift for each value gives a continuous
lift. The two other exactness assertions hold pointwise, with the
induced topology on each subgroup still discrete. This gives a
short exact sequence of cochain complexes. Its connecting arrow
is obtained by lifting a cocycle and applying the differential,
as in \emph{Algebra Core}, Theorem 10.3
\cite[PDF p.~51]{mg-algebra-core}.
\end{proof}

\begin{theorem}[Kummer isomorphism with its Galois action]
\label{res:GH06}
Let \(m\geq 1\) and assume that \(m\) is invertible in \(k\). Give
\(\mu_m\subset(k^{\mathrm s})^\times\) its natural \(G_k\)-action.
Then there is an isomorphism
\[
 k^\times/k^{\times m}
   \ \xrightarrow{\ \sim\ }\ H^1_{\mathrm{cts}}(G_k,\mu_m),
 \qquad
 a\longmapsto [\,g\mapsto g(\alpha)/\alpha\,],\quad \alpha^m=a.
\]
If \(\mu_m\subset k^\times\), this identifies the quotient with
\(\operatorname{Hom}_{\mathrm{cts}}(G_k,\mu_m)\).
\end{theorem}
\begin{proof}
The \(m\)-th power on \((k^{\mathrm s})^\times\) is surjective.
Indeed \(X^m-u\) is separable for \(u\ne0\), by its derivative, and
splits over the separably closed \(k^{\mathrm s}\). Its kernel on
units is \(\mu_m\). The maps are equivariant, and all coefficients
are continuous by the stabilizers in the Krull topology.
Invariants in \((k^{\mathrm s})^\times\) are \(k^\times\), by
Result C02 of \emph{Infinite Galois Theory}
\cite[PDF p.~12]{mg-aft}. Proposition~\ref{res:GH05} gives the connecting
map of the short exact sequence, and Theorem~\ref{res:GH03} makes it
surjective onto \(H^1_{\mathrm{cts}}(G_k,\mu_m)\). Its formula follows
from the coboundary of \(\alpha\). Changing \(\alpha\) by an element of
\(\mu_m\) multiplies the cocycle by that element's coboundary. For
trivial action the cocycle identity is the homomorphism condition and
every coboundary is the identity.
\end{proof}

\begin{theorem}[Finite Kummer duality]
\label{res:GH07}
Assume that \(m\) is invertible in \(k\) and
\(\mu_m\subset k^\times\). Inside \(k^{\mathrm s}\), the finite Galois
extensions with abelian group of exponent dividing \(m\)
correspond to the finite subgroups \(D\subset k^\times/k^{\times m}\).
The field \(L_D\) is obtained by adjoining \(m\)-th roots of
representatives of \(D\). If \(L\) is one of the extensions, its
subgroup is
\[
 D_L=\bigl(k^\times\cap (L^\times)^m\bigr)/k^{\times m}.
\]
The pairing \((g,[a])\mapsto g(\alpha)/\alpha\), with \(\alpha^m=a\),
identifies either of \(\operatorname{Gal}(L_D/k)\) and \(D\) with the
group of homomorphisms from the other to \(\mu_m\).
\end{theorem}
\begin{proof}
Any finite subgroup of the units of a field is cyclic. Its exponent
\(e\) is realized as the order of an element: choose the prime-power
parts of elements having maximal such parts and multiply them.
Then the polynomial \(X^e-1\) bounds the order of the subgroup by \(e\).
In particular \(\mu_m\) is cyclic of order \(m\), since \(X^m-1\) is
separable and split in \(k^{\mathrm s}\).

For a finite abelian group \(A\) of exponent dividing \(m\), its
\(\mu_m\)-valued characters separate elements and number \(|A|\).
Indeed, a character on a subgroup \(H\) extends across an element
\(x\) of order \(d\) modulo \(H\). Since \(d\mid m\),
\(\chi(x^d)^{m/d}=1\); choose \(\zeta\in\mu_m\) with
\(\zeta^d=\chi(x^d)\). Then
\(\chi(hx^i)=\chi(h)\zeta^i\) gives the extension. Repeating proves
surjectivity of restriction to any subgroup. This extends a faithful
character of each cyclic subgroup. Counting the characters by
restriction to a nontrivial cyclic subgroup and using the quotient
for the kernel gives the count by induction on \(|A|\).

Choose finitely many generators of \(D\). The field generated by
each chosen root is finite cyclic Galois by the single-radical
result \texttt{K07.15.kummer\_forward} in D07
\cite[PDF p.~135]{mg-d07}. By Result F04 of
\emph{Infinite Galois Theory} \cite[PDF p.~5]{mg-aft},
their compositum \(L_D/k\) is finite Galois.
If \([a]\) is a product of powers of the generator classes, an
\(m\)-th root of \(a\) in \(L_D\) is the corresponding product of
powers of the chosen roots times a scalar in \(k^\times\).
On these generators the ratios give an injective homomorphism
of \(\operatorname{Gal}(L_D/k)\) into a product of \(\mu_m\)'s.
The pairing is independent of roots and representatives and is
multiplicative in both variables. Its right kernel is trivial by
the fixed-field theorem: if each \(g\) fixes \(\alpha\), then
\(\alpha\in k\), so \([a]=1\). The left kernel is trivial because
the chosen roots generate \(L_D\). The character counts just proved
now make both maps isomorphisms.

For a finite Galois \(L/k\) with abelian group \(A\) of exponent
dividing \(m\), each character \(\chi:A\to\mu_m\) is a cocycle on
\(L^\times\). Theorem~\ref{res:GH03} gives \(\chi(g)=g(b)/b\), with
\(b\in L^\times\) and \(b^m\in k^\times\). Choosing such a \(b\) for
each character generates \(L\): the subgroup fixing them all is
trivial because the characters separate, and the finite Galois
correspondence applies (Result F01 of \emph{Infinite Galois Theory}
\cite[PDF p.~4]{mg-aft}). Finally \([a]\mapsto(g\mapsto g(\alpha)/\alpha)\)
injects \(D_L\) into the character group, by the fixed-field theorem,
and the chosen elements show surjectivity. For \(L=L_D\), inclusion
\(D\subseteq D_L\) and the equal sizes from perfectness give equality.
\end{proof}

\section{Artin--Schreier theory}

\begin{proposition}[Trace contraction for semilinear coefficients]
\label{res:GH08}
Let \(E/k\) be finite Galois with group \(H\), and let \(V\) be an
\(E\)-vector space on which \(H\) acts semilinearly, so
\(\sigma(av)=\sigma(a)\sigma(v)\). Its additive group has
\(H^q(H,V^+)=0\) for all \(q>0\). For any Galois extension \(L/k\),
\[
 H^q_{\mathrm{cts}}\bigl(\operatorname{Gal}(L/k),L^+\bigr)=0
 \qquad(q>0).
\]
\end{proposition}
\begin{proof}
Use equivariant homogeneous cochains \(f:H^{q+1}\to V^+\),
with differential the alternating sum of omissions.
The change to inhomogeneous cochains evaluates \(f\) at
\((1,g_1,g_1g_2,\ldots,g_1\cdots g_q)\); its inverse takes
\(F\) to \(g_0 F(g_0^{-1}g_1,\ldots,g_{q-1}^{-1}g_q)\).
Deleting the first variable yields the action term and deleting
an interior variable multiplies adjacent increments. These formulas
intertwine the two differentials.
Choose trace-one \(t\) by Lemma~\ref{res:GH02}. For \(q\geq1\) put
\[
 (s f)(g_0,\ldots,g_{q-1})
   =\sum_{h\in H}h(t)\,f(h,g_0,\ldots,g_{q-1}).
\]
The reindexing \(h=\sigma h'\) and semilinearity prove
equivariance. In \(s d f\), omitting \(h\) gives
\(\sum_h h(t)f=f\); each of the remaining terms cancels its
counterpart in \(d s f\) by the shift of signs. Hence \(d s+s d=1\)
on positive degrees. For \(L/k\) the fixed coefficients at stage
\(E\) are \(E^+\), so use Proposition~\ref{res:GH01}.
\end{proof}

\begin{theorem}[The Artin--Schreier connecting isomorphism]
\label{res:GH09}
If \(\operatorname{char}k=p>0\) and \(\wp(x)=x^p-x\), then
\[
 k/\wp(k)
   \ \cong\ H^1_{\mathrm{cts}}(G_k,\mathbb F_p)
   =\operatorname{Hom}_{\mathrm{cts}}(G_k,\mathbb F_p),
 \qquad
 [a]\longmapsto \bigl(g\mapsto g(\alpha)-\alpha\bigr)
 \quad(\wp(\alpha)=a).
\]
The coefficient action on \(\mathbb F_p\) is trivial.
\end{theorem}
\begin{proof}
The polynomial \(X^p-X-u\) is separable, since its derivative is
\(-1\), so \(\wp:k^{\mathrm s}\to k^{\mathrm s}\) is surjective.
Its kernel is exactly \(\mathbb F_p\), which already gives \(p\)
roots. The resulting sequence of discrete coefficients is exact.
Apply Proposition~\ref{res:GH05}, the invariants formula
(Result C02 of \emph{Infinite Galois Theory}
\cite[PDF p.~12]{mg-aft}), and Proposition~\ref{res:GH08} in
degree one. The connecting cocycle is the difference displayed.
\end{proof}

For a nonzero class its character has kernel
\(\operatorname{Gal}(k^{\mathrm s}/k(\alpha))\), since the value at
\(g\) is zero exactly when \(g\) fixes \(\alpha\). This realizes
the cyclic degree-\(p\) extensions of D07,
\texttt{K07.20.artin\_schreier} \cite[PDF p.~172]{mg-d07},
as the lines in the character space: changing the isomorphism from
the cyclic quotient to \(\mathbb F_p^+\) rescales the character.

\begin{theorem}[Vanishing for \(p\)-primary coefficients in characteristic \(p\)]
\label{res:GH10}
In characteristic \(p>0\), every discrete continuous \(G_k\)-module
\(M\) with \(p\)-primary torsion underlying group satisfies
\[
 H^q_{\mathrm{cts}}(G_k,M)=0\qquad(q\geq2).
\]
\end{theorem}
\begin{proof}
First suppose \(pM=0\), and put
\(W=k^{\mathrm s}\otimes_{\mathbb F_p}M\), with diagonal action and
the discrete topology. It is continuous because each tensor is a
finite sum of elementary tensors, and each factor has an open
stabilizer. The Artin--Schreier sequence tensored over \(\mathbb F_p\)
with \(M\) gives
\[
 0\longrightarrow M\xrightarrow{m\mapsto 1\otimes m}W
 \xrightarrow{\wp\otimes 1}W\longrightarrow0.
\]
Exactness follows by choosing a vector-space basis of \(M\), which
turns the tensor sequence into a direct sum of copies of the
Artin--Schreier sequence from Theorem~\ref{res:GH09}. The maps in the
display are equivariant because \(\wp\) commutes with field
automorphisms.

For any continuous \(q\)-cocycle with values in \(W\), finitely
many field elements and vectors occur in expressions for its
finitely many values. A vector in the discrete \(M\) has finite
orbit by compactness of \(G_k\). Thus there is a finite-dimensional
stable \(\mathbb F_p\)-subspace \(M_0\) containing those vectors.
Choose a finite Galois layer \(E\) containing the field elements,
large enough that \(U_E\) fixes \(M_0\) and the cocycle factors
through \(G_E^q\). For the latter use the cochain factorization in
the proof of Proposition~\ref{res:GH01}. For the former use the
intersection of the stabilizers of a finite basis. The cocycle
then takes values in \(E\otimes_{\mathbb F_p} M_0\), on which
\(G_E\) acts semilinearly. This group embeds in \(W\) since tensors
over a field preserve injections, and restriction onto \(G_E\)
is surjective. The cocycle identity holds at that stage. For
\(q>0\), Proposition~\ref{res:GH08} kills the class at that stage,
so \(H^q_{\mathrm{cts}}(G_k,W)=0\). The sequence above and
Proposition~\ref{res:GH05} prove the assertion for \(pM=0\).

For \(M\) killed by \(p^r\), use the exact sequence
\(0\to M[p]\to M\xrightarrow{p}pM\to0\) and induction on \(r\).
For general \(p\)-primary torsion \(M\), each continuous cocycle
has finite image and takes values in \(M[p^r]\) for some \(r\).
Its coboundary primitive found in that submodule also works in \(M\).
\end{proof}

\section{Cyclic \texorpdfstring{\(H^2\)}{H2} and arithmetic calculations}

\begin{proposition}[Second cohomology of a finite cyclic group]
\label{res:GH11}
Let \(C=\langle s\rangle\) have order \(d\geq1\) and \(A\) be an
abelian \(C\)-module. For \(N=1+s+\cdots+s^{d-1}\), one has
\[
 H^2(C,A)\cong A^C/N(A).
\]
In particular for a finite cyclic field extension \(E/k\),
\[
 H^2\bigl(\operatorname{Gal}(E/k),E^\times\bigr)
     \cong k^\times/N_{E/k}(E^\times).
\]
\end{proposition}
\begin{proof}
The order-one case is immediate from the one-element group complex.
For \(d\geq2\) use additive notation. Every class has a normalized
cocycle \(c\), with \(c(1,g)=c(g,1)=0\). In fact the cocycle identity
implies \(c(1,g)=c(1,1)\) and \(c(g,1)=g c(1,1)\); subtract the
coboundary of the constant one-cochain \(c(1,1)\).
Define a group on \(A\times C\) by
\((a,g)(b,h)=(a+g b+c(g,h),gh)\).
Associativity is the cocycle identity, the normalized values give
the identity \((0,1)\), and inverse elements are obtained by first
solving the equation in \(C\) and then in \(A\).
The element \(u=(0,s)\) conjugates \(A\) by \(s\), and
\(u^d=(a,1)\) for some \(a\in A^C\), because \(u\) commutes with
\(u^d\). A different lift \(b u\), for \(b\in A\), replaces
\(a\) by \(a+N b\) after moving the \(d\) copies of \(b\) to the
left. Replacing \(c\) by \(c+df\) induces the isomorphism
\((b,g)\mapsto (b+f(g),g)\) from the new group to the old one, so the
class of \(a\) is well defined.

Use the section \(s^i\mapsto u^i\) for \(0\leq i<d\).
Changing a section changes its factor set by a coboundary:
if each section element is multiplied on the left by \(f(g)\in A\),
its factor set at \((g,h)\) changes by
\(f(g)+g f(h)-f(gh)\). The factor set of this section has value \(0\)
at \((s^i,s^j)\) when \(i+j<d\) and \(a\) otherwise.
Conversely, for every \(a\in A^C\) this carry rule defines a
cocycle: adding \(i,j,\ell\) and reducing modulo \(d\) in either
order gives the same total carry because \(a\) is invariant.
Its \(d\)-th power computation gives \(a\), and its cohomology
class depends additively on \(a\). If two such cocycles have
parameters differing by \(N b\), the section defined by the
new lift \(b u\) changes one into the other. These checks give the
isomorphism. The field formula follows by the finite fixed-field
theorem and the multiplicative norm.
\end{proof}

\begin{example}[Finite fields]
\label{res:GH12}
For \(k=\mathbb F_q\) of characteristic \(p\) and \(p\nmid m\),
\(H^1_{\mathrm{cts}}(G_k,\mu_m)\) is cyclic of order
\(\gcd(m,q-1)\). The Artin--Schreier quotient \(k/\wp(k)\)
is identified with \(\mathbb F_p\) by the trace, and the
character of \(a\) takes the trace value at the arithmetic
Frobenius \(x\mapsto x^q\). Also
\[
 H^2_{\mathrm{cts}}\bigl(G_k,(k^{\mathrm s})^\times\bigr)=0.
\]
\end{example}
\begin{proof}
Apply Theorem~\ref{res:GH06} and cyclicity of the units, proved in
Theorem~\ref{res:GH07}. Write \(q=p^r\). The trace
\(\operatorname{Tr}_{\mathbb F_q/\mathbb F_p}\) is
\(\sum_{i=0}^{r-1}a^{p^i}\), by the Frobenius description of
finite Galois layers (Results E01--E04 of \emph{Infinite Galois
Theory} \cite[PDF pp.~18--19]{mg-aft}), and induces the claimed
quotient by \texttt{K07.20.finite\_artin\_schreier} of D07
\cite[PDF p.~173]{mg-d07}. For a root \(\wp(\alpha)=a\),
the equality \(\alpha^q-\alpha=\sum_{i=0}^{r-1}a^{p^i}\) follows
by raising the defining equation to \(p^i\) and summing. This gives
the stated evaluation in Theorem~\ref{res:GH09}.
If \(E=\mathbb F_{q^d}\), its norm on units is the power
\(1+q+\dots+q^{d-1}=(q^d-1)/(q-1)\). A generator therefore has norm
of order \(q-1\), so the norm is onto. The finite-layer
\(H^2(G_E,E^\times)\) vanishes by Proposition~\ref{res:GH11}.
Proposition~\ref{res:GH01} now gives the last equality.
\end{proof}

\begin{example}[A real norm obstruction and two rational-function classes]
\label{res:GH13}
For \(\mathbb C/\mathbb R\), the relative group
\(H^2(\operatorname{Gal}(\mathbb C/\mathbb R),\mathbb C^\times)\)
has order two, with nonzero class represented by a cyclic cocycle
carrying the value \(-1\).
Over \(k=\mathbb F_q(t)\) of characteristic \(p\), \(a=t\) gives a
nonzero Artin--Schreier class. If \(m>1\) divides \(q-1\), the Kummer
class of \(t\) has order \(m\) and gives a cyclic extension of degree
\(m\).
\end{example}
\begin{proof}
The complex norm has image exactly the positive real numbers, by
\(z\bar z>0\) for \(z\ne0\) and positive real square roots. The
first assertion is then Proposition~\ref{res:GH11} with multiplicative
notation. For \(t\), put \(v_\infty(f/g)=\deg g-\deg f\) on
nonzero rational functions in coprime polynomials, and
\(v_\infty(0)=+\infty\). A common denominator shows
\(v_\infty(f+g)\geq\min(v_\infty(f),v_\infty(g))\), with equality
for unequal values because the leading terms cannot cancel.
If \(v_\infty(x)\geq0\), then \(v_\infty(\wp(x))\geq0\); if
\(v_\infty(x)<0\), then \(v_\infty(\wp(x))=p v_\infty(x)\).
Neither gives the value \(-1\) of \(t\), so
Theorem~\ref{res:GH09} applies. Finally \(v_\infty(t^i)=-i\)
and the valuation of any \(m\)-th power is divisible by \(m\).
The class of \(t\) has order \(m\), and \(m\mid q-1\) puts
\(\mu_m\) inside \(k\) by cyclicity of \(\mathbb F_q^\times\)
from Theorem~\ref{res:GH07}. The perfect pairing there gives the
cyclic degree-\(m\) extension.
\end{proof}

\section{Inflation and restriction}

\begin{proposition}[Continuous inflation--restriction in degree one]
\label{res:GH14}
For a profinite group \(G\), a closed normal subgroup \(N\), and a
discrete continuous abelian \(G\)-module \(M\), the sequence
\[
 0\longrightarrow H^1_{\mathrm{cts}}(G/N,M^N)
 \xrightarrow{\mathrm{inf}} H^1_{\mathrm{cts}}(G,M)
 \xrightarrow{\mathrm{res}} H^1_{\mathrm{cts}}(N,M)^{G/N}
\]
is exact at its first two cohomology groups. The action on the last
group comes from \((g\cdot f)(n)=g f(g^{-1}ng)\).
\end{proposition}
\begin{proof}
The displayed action respects the continuous cocycle identity
and sends the coboundary of \(b\in M\) to that of \(gb\).
For any cocycle \(c\) on \(G\), expanding the three factors in
\(g^{-1}ng\) gives
\[
 g c(g^{-1}ng)=c(n)+(n-1)c(g).
\]
The same identity within \(N\), for any cocycle on \(N\), shows
that inner action by elements of \(N\) is trivial on cohomology.
Thus the action factors through \(G/N\), and the identity for
\(c\) on \(G\) places its restriction class among the invariants.

An inflated cocycle vanishes on \(N\). If it is a coboundary
\(g\mapsto(g-1)b\) on \(G\), then \(b\in M^N\), so its original
class was already zero. Conversely a class restricting to zero
has a representative \(c\) vanishing on \(N\), obtained by subtracting
a coboundary whose restriction is the given restriction cocycle.
The cocycle rule gives \(c(gn)=c(g)\) and \(c(ng)=n c(g)\).
Normality gives \(ng\in gN\), so \(n c(g)=c(g)\). Hence \(c\)
descends to a cocycle on \(G/N\) with values in \(M^N\).
It is continuous by the quotient topology. The action of \(G/N\)
on the discrete \(M^N\) is continuous because its point stabilizers
are the images of the open stabilizers in \(G\), each containing
\(N\). Closed normal quotients here are profinite by Result PF02
of \emph{Topics in Group Theory} \cite[PDF pp.~15--16]{mg-groups}.
\end{proof}

\paragraph{Sources.}
The finite-layer and degree-one arithmetic constructions are in
\cite[Sections 2.2, 2.4--2.5]{sharifi} and
\cite[Chapter~I, Sections 2--3 and 5]{serregc}.
For field-theoretic independence, norm and radical expressions,
see \cite[Chapter~5]{milneft}.
